% GENERATED FILE. Edit canonical lessons, then run npm run generate:books.
% canonical-source-hash: fnv1a64:43445a541c74eacf
% canonical-lessons: PT-C524-poesia, PT-C524-poeta, PT-C524-verso, PT-C524-rima, PT-C524-declamar

\chapter{Versos e rimas}
\label{ch:pt-versos-e-rimas}
\hlchaptermodality{524}

\begin{chapteropening}
\textbf{By the end of this chapter:} \emph{I can talk about poetry.}

Complete the last lesson of chapter 524: I can talk about poetry.
\end{chapteropening}

\section[a poesia]{a poesia — poetry; poem}

\label{lesson:PT-C524-poesia}



\begin{quote}
Before the new one: say the Portuguese for to reread, then the Portuguese for bookmark.
\end{quote}

\subsection*{You'll want to know: a poesia}
\textbf{a poesia} — \textquotedblleft{}poetry; poem\textquotedblright{}.

\begin{quote}
A minha avó sabe muita poesia de cor.
\end{quote}

\begin{quote}
\emph{My grandmother knows a lot of poetry by heart.}
\end{quote}

\emph{Uma poesia} can also mean one poem.

\begin{cousinweb}
\emph{poesia} comes from Latin \emph{poēsis} (EtymDB), from Greek \emph{poíēsis} \textquotedblleft{}a making, a composing\textquotedblright{}, built on \emph{poieîn} \textquotedblleft{}to make\textquotedblright{}.
\end{cousinweb}

\begin{grammarlens}[title={What you've built}]
One more word for telling, explaining and describing at B1.
\end{grammarlens}

\subsection*{Your turn}
\begin{itemize}
\raggedright
  \item \emph{Say it:} \emph{a poesia}
  \item \emph{Say it:} \emph{a poesia}, once more
  \item \emph{Say it:} a sentence of your own with \emph{a poesia}
  \item \emph{Recall:} say the Portuguese for to reread, then the Portuguese for bookmark, then say \emph{a poesia} again
\end{itemize}

\begin{tcolorbox}[breakable,colback=teal!4,colframe=teal!35!black,title={Before you move on}]
What does a poesia mean? (\textquotedblleft{}Poetry; poem\textquotedblright{}.) Say it once more.
\end{tcolorbox}
\section[o poeta]{o poeta — poet}

\label{lesson:PT-C524-poeta}



\begin{quote}
Before the new one: say the Portuguese for bookmark, then the Portuguese for poetry.
\end{quote}

\subsection*{You'll want to know: o poeta}
\textbf{o poeta} — \textquotedblleft{}poet\textquotedblright{}.

\begin{quote}
O poeta leu os seus poemas na biblioteca.
\end{quote}

\begin{quote}
\emph{The poet read his poems in the library.}
\end{quote}

A woman is \emph{a poetisa} or \emph{a poeta}.

\begin{cousinweb}
\emph{poeta} comes from Latin \emph{poēta} \textquotedblleft{}a poet\textquotedblright{} (EtymDB), from Greek \emph{poiēt\'{\={e}}s} \textquotedblleft{}a maker\textquotedblright{}, built on \emph{poieîn} \textquotedblleft{}to make\textquotedblright{}.
\end{cousinweb}

\begin{grammarlens}[title={What you've built}]
One more word for telling, explaining and describing at B1.
\end{grammarlens}

\subsection*{Your turn}
\begin{itemize}
\raggedright
  \item \emph{Say it:} \emph{o poeta}
  \item \emph{Say it:} \emph{o poeta}, once more
  \item \emph{Say it:} a sentence of your own with \emph{o poeta}
  \item \emph{Recall:} say the Portuguese for bookmark, then the Portuguese for poetry, then say \emph{o poeta} again
\end{itemize}

\begin{tcolorbox}[breakable,colback=teal!4,colframe=teal!35!black,title={Before you move on}]
What does o poeta mean? (\textquotedblleft{}Poet\textquotedblright{}.) Say it once more.
\end{tcolorbox}
\section[o verso]{o verso — line (of poetry), verse; back (of a page)}

\label{lesson:PT-C524-verso}



\begin{quote}
Before the new one: say the Portuguese for poetry, then the Portuguese for poet.
\end{quote}

\subsection*{You'll want to know: o verso}
\textbf{o verso} — \textquotedblleft{}line (of poetry), verse; back (of a page)\textquotedblright{}.

\begin{quote}
Este verso é o meu preferido.
\end{quote}

\begin{quote}
\emph{This line is my favourite.}
\end{quote}

Also the back of a sheet: \emph{no verso da folha}.

\begin{cousinweb}
\emph{verso} is inherited from Latin \emph{versus} \textquotedblleft{}a turning, a furrow, a line of writing\textquotedblright{} (EtymDB), from \emph{vertere} \textquotedblleft{}to turn\textquotedblright{}.
\end{cousinweb}

\begin{grammarlens}[title={What you've built}]
One more word for telling, explaining and describing at B1.
\end{grammarlens}

\subsection*{Your turn}
\begin{itemize}
\raggedright
  \item \emph{Say it:} \emph{o verso}
  \item \emph{Say it:} \emph{o verso}, once more
  \item \emph{Say it:} a sentence of your own with \emph{o verso}
  \item \emph{Recall:} say the Portuguese for poetry, then the Portuguese for poet, then say \emph{o verso} again
\end{itemize}

\begin{tcolorbox}[breakable,colback=teal!4,colframe=teal!35!black,title={Before you move on}]
What does o verso mean? (\textquotedblleft{}Line (of poetry), verse; back (of a page)\textquotedblright{}.) Say it once more.
\end{tcolorbox}
\section[a rima]{a rima — rhyme}

\label{lesson:PT-C524-rima}



\begin{quote}
Before the new one: say the Portuguese for poet, then the Portuguese for line.
\end{quote}

\subsection*{You'll want to know: a rima}
\textbf{a rima} — \textquotedblleft{}rhyme\textquotedblright{}.

\begin{quote}
As canções infantis têm muitas rimas.
\end{quote}

\begin{quote}
\emph{Children's songs have lots of rhymes.}
\end{quote}

The verb is \emph{rimar}: \emph{amor rima com dor}.

\begin{cousinweb}
\emph{rima} goes back to Medieval Latin \emph{rithmus} (EtymDB), which is usually traced to Greek \emph{rhythmós} \textquotedblleft{}measured flow, rhythm\textquotedblright{}.
\end{cousinweb}

\begin{grammarlens}[title={What you've built}]
One more word for telling, explaining and describing at B1.
\end{grammarlens}

\subsection*{Your turn}
\begin{itemize}
\raggedright
  \item \emph{Say it:} \emph{a rima}
  \item \emph{Say it:} \emph{a rima}, once more
  \item \emph{Say it:} a sentence of your own with \emph{a rima}
  \item \emph{Recall:} say the Portuguese for poet, then the Portuguese for line, then say \emph{a rima} again
\end{itemize}

\begin{tcolorbox}[breakable,colback=teal!4,colframe=teal!35!black,title={Before you move on}]
What does a rima mean? (\textquotedblleft{}Rhyme\textquotedblright{}.) Say it once more.
\end{tcolorbox}
\section[declamar]{declamar — to recite (aloud)}

\label{lesson:PT-C524-declamar}



\begin{quote}
Before the new one: say the Portuguese for line, then the Portuguese for rhyme.
\end{quote}

\subsection*{You'll want to know: declamar}
\textbf{declamar} — \textquotedblleft{}to recite (aloud)\textquotedblright{}.

\begin{quote}
No fim da festa, as crianças declamaram poemas.
\end{quote}

\begin{quote}
\emph{At the end of the party, the children recited poems.}
\end{quote}

More solemn than \emph{ler em voz alta}.

\begin{cousinweb}
\emph{declamar} is a learned borrowing from Latin \emph{dēclāmāre} \textquotedblleft{}to practise speaking aloud\textquotedblright{} (Lewis \& Short), built on \emph{clāmāre} \textquotedblleft{}to shout\textquotedblright{}.
\end{cousinweb}

\begin{grammarlens}[title={What you've built}]
One more word for telling, explaining and describing at B1.
\end{grammarlens}

\subsection*{Your turn}
\begin{itemize}
\raggedright
  \item \emph{Say it:} \emph{declamar}
  \item \emph{Say it:} \emph{declamar}, once more
  \item \emph{Say it:} a sentence of your own with \emph{declamar}
  \item \emph{Recall:} say the Portuguese for line, then the Portuguese for rhyme, then say \emph{declamar} again
\end{itemize}

\begin{tcolorbox}[breakable,colback=teal!4,colframe=teal!35!black,title={Before you move on}]
What does declamar mean? (\textquotedblleft{}To recite (aloud)\textquotedblright{}.) Say it once more.
\end{tcolorbox}
